% GENERATED FILE. Edit canonical lessons, then run npm run generate:books.
% canonical-source-hash: fnv1a64:62ef87aeda8f31ef
% canonical-lessons: GU-C369-bhay, GU-C369-dahesat, GU-C369-nirat, GU-C369-asanka, GU-C369-daravvu

\chapter{Fears}
\label{ch:gu-fears}
\hlchaptermodality{369}

\begin{chapteropening}
\textbf{By the end of this chapter:} \emph{I can say what I am afraid might happen.}

Complete the last lesson of chapter 369: I can say what I am afraid might happen.
\end{chapteropening}

\section[\texorpdfstring{\gu{ભય} (bhay)}{ભય (bhay)}]{\gu{ભય} (bhay) — dread}

\label{lesson:GU-C369-bhay}



\begin{quote}
Before the new one: say the Gujarati for to bear fruit, then the Gujarati for to turn out.
\end{quote}

\subsection*{You'll want to know: \gu{ભય}}
\textbf{\gu{ભય}} — \emph{bhay} — \textquotedblleft{}dread\textquotedblright{}.

\begin{quote}
\gu{જો} \gu{નદીમાં} \gu{પૂર} \gu{આવશે} \gu{તો} \gu{ગામને} \gu{ભય} \gu{છે}.
\end{quote}

\begin{quote}
\emph{jo nadīmā̃ pūr āvśe to gāmne bhay che.}
\end{quote}

\begin{quote}
If the river floods, the village is in danger.
\end{quote}

A masculine noun.

\begin{cousinweb}
\textbf{\gu{ભય}} (\emph{bhay}) is a tatsama, borrowed from Sanskrit \emph{bhaya}, \textquotedblleft{}fear, alarm, dread\textquotedblright{} (Monier-Williams).
\end{cousinweb}

\begin{grammarlens}[title={What you've built}]
One more word for saying what happens if.
\end{grammarlens}

\subsection*{Your turn}
\begin{itemize}
\raggedright
  \item \emph{Say it:} \emph{bhay}
  \item \emph{Say it:} \emph{bhay}, once more
  \item \emph{Say it:} \emph{jo nadīmā̃ pūr āvśe to gāmne bhay che.} — If the river floods, the village is in danger.
  \item \emph{Recall:} say the Gujarati for to bear fruit, then the Gujarati for to turn out, then say \emph{bhay} again
\end{itemize}

\begin{tcolorbox}[breakable,colback=teal!4,colframe=teal!35!black,title={Before you move on}]
What does \gu{ભય} mean? (\textquotedblleft{}Dread\textquotedblright{}.) Say it once more.
\end{tcolorbox}
\section[\texorpdfstring{\gu{દહેશત} (daheśat)}{દહેશત (daheśat)}]{\gu{દહેશત} (daheśat) — terror}

\label{lesson:GU-C369-dahesat}



\begin{quote}
Before the new one: say the Gujarati for to turn out, then the Gujarati for dread.
\end{quote}

\subsection*{You'll want to know: \gu{દહેશત}}
\textbf{\gu{દહેશત}} — \emph{daheśat} — \textquotedblleft{}terror\textquotedblright{}.

\begin{quote}
\gu{જો} \gu{વાઘ} \gu{જંગલમાંથી} \gu{બહાર} \gu{આવશે} \gu{તો} \gu{ગામમાં} \gu{દહેશત} \gu{ફેલાશે}.
\end{quote}

\begin{quote}
\emph{jo vāgh jangalmā̃thī bahār āvśe to gāmmā̃ daheśat phelāśe.}
\end{quote}

\begin{quote}
If the tiger comes out of the forest, terror will spread through the village.
\end{quote}

A feminine noun.

\begin{cousinweb}
\textbf{\gu{દહેશત}} (\emph{daheśat}) is usually traced to Persian \emph{dahshat}, \textquotedblleft{}terror, amazement\textquotedblright{}, a word of Arabic origin.
\end{cousinweb}

\begin{grammarlens}[title={What you've built}]
One more word for saying what happens if.
\end{grammarlens}

\subsection*{Your turn}
\begin{itemize}
\raggedright
  \item \emph{Say it:} \emph{daheśat}
  \item \emph{Say it:} \emph{daheśat}, once more
  \item \emph{Say it:} \emph{jo vāgh jangalmā̃thī bahār āvśe to gāmmā̃ daheśat phelāśe.} — If the tiger comes out of the forest, terror will spread through the village.
  \item \emph{Recall:} say the Gujarati for to turn out, then the Gujarati for dread, then say \emph{daheśat} again
\end{itemize}

\begin{tcolorbox}[breakable,colback=teal!4,colframe=teal!35!black,title={Before you move on}]
What does \gu{દહેશત} mean? (\textquotedblleft{}Terror\textquotedblright{}.) Say it once more.
\end{tcolorbox}
\section[\texorpdfstring{\gu{નિરાંત} (nirā̃t)}{નિરાંત (nirā̃t)}]{\gu{નિરાંત} (nirā̃t) — peace of mind; relief}

\label{lesson:GU-C369-nirat}



\begin{quote}
Before the new one: say the Gujarati for dread, then the Gujarati for terror.
\end{quote}

\subsection*{You'll want to know: \gu{નિરાંત}}
\textbf{\gu{નિરાંત}} — \emph{nirā̃t} — \textquotedblleft{}peace of mind; relief\textquotedblright{}.

\begin{quote}
\gu{જો} \gu{પરિણામ} \gu{સારું} \gu{આવશે} \gu{તો} \gu{મને} \gu{નિરાંત} \gu{થશે}.
\end{quote}

\begin{quote}
\emph{jo pariṇām sārũ āvśe to mane nirā̃t thaśe.}
\end{quote}

\begin{quote}
If the result is good, I will be relieved.
\end{quote}

A feminine noun.

\begin{cousinweb}
\textbf{\gu{નિરાંત}} (\emph{nirā̃t}) CDIAL lists the Gujarati noun, \textquotedblleft{}peace, tranquillity\textquotedblright{}, under a reconstructed \emph{nirānti} (CDIAL 7273).
\end{cousinweb}

\begin{grammarlens}[title={What you've built}]
One more word for saying what happens if.
\end{grammarlens}

\subsection*{Your turn}
\begin{itemize}
\raggedright
  \item \emph{Say it:} \emph{nirā̃t}
  \item \emph{Say it:} \emph{nirā̃t}, once more
  \item \emph{Say it:} \emph{jo pariṇām sārũ āvśe to mane nirā̃t thaśe.} — If the result is good, I will be relieved.
  \item \emph{Recall:} say the Gujarati for dread, then the Gujarati for terror, then say \emph{nirā̃t} again
\end{itemize}

\begin{tcolorbox}[breakable,colback=teal!4,colframe=teal!35!black,title={Before you move on}]
What does \gu{નિરાંત} mean? (\textquotedblleft{}Peace of mind; relief\textquotedblright{}.) Say it once more.
\end{tcolorbox}
\section[\texorpdfstring{\gu{આશંકા} (āśankā)}{આશંકા (āśankā)}]{\gu{આશંકા} (āśankā) — a misgiving}

\label{lesson:GU-C369-asanka}



\begin{quote}
Before the new one: say the Gujarati for terror, then the Gujarati for peace of mind.
\end{quote}

\subsection*{You'll want to know: \gu{આશંકા}}
\textbf{\gu{આશંકા}} — \emph{āśankā} — \textquotedblleft{}a misgiving\textquotedblright{}.

\begin{quote}
\gu{મને} \gu{આશંકા} \gu{છે} \gu{કે} \gu{જો} \gu{આપણે} \gu{મોડા} \gu{પડીશું} \gu{તો} \gu{ટ્રેન} \gu{છૂટી} \gu{જશે}.
\end{quote}

\begin{quote}
\emph{mane āśankā che ke jo āpaṇe moḍā paḍīśũ to ṭren chūṭī jaśe.}
\end{quote}

\begin{quote}
I am afraid that if we are late we will miss the train.
\end{quote}

A feminine noun.

\begin{cousinweb}
\textbf{\gu{આશંકા}} (\emph{āśankā}) is a tatsama, borrowed from Sanskrit \emph{āśaṅkā}, \textquotedblleft{}fear, apprehension\textquotedblright{} (Monier-Williams).
\end{cousinweb}

\begin{grammarlens}[title={What you've built}]
One more word for saying what happens if.
\end{grammarlens}

\subsection*{Your turn}
\begin{itemize}
\raggedright
  \item \emph{Say it:} \emph{āśankā}
  \item \emph{Say it:} \emph{āśankā}, once more
  \item \emph{Say it:} \emph{mane āśankā che ke jo āpaṇe moḍā paḍīśũ to ṭren chūṭī jaśe.} — I am afraid that if we are late we will miss the train.
  \item \emph{Recall:} say the Gujarati for terror, then the Gujarati for peace of mind, then say \emph{āśankā} again
\end{itemize}

\begin{tcolorbox}[breakable,colback=teal!4,colframe=teal!35!black,title={Before you move on}]
What does \gu{આશંકા} mean? (\textquotedblleft{}A misgiving\textquotedblright{}.) Say it once more.
\end{tcolorbox}
\section[\texorpdfstring{\gu{ડરાવવું} (ḍarāvvũ)}{ડરાવવું (ḍarāvvũ)}]{\gu{ડરાવવું} (ḍarāvvũ) — to frighten}

\label{lesson:GU-C369-daravvu}



\begin{quote}
Before the new one: say the Gujarati for peace of mind, then the Gujarati for a misgiving.
\end{quote}

\subsection*{You'll want to know: \gu{ડરાવવું}}
\textbf{\gu{ડરાવવું}} — \emph{ḍarāvvũ} — \textquotedblleft{}to frighten\textquotedblright{}.

\begin{quote}
\gu{જો} \gu{તમે} \gu{બાળકને} \gu{ડરાવશો} \gu{તો} \gu{તે} \gu{રાત્રે} \gu{નહીં} \gu{સૂએ}.
\end{quote}

\begin{quote}
\emph{jo tame bāḷakne ḍarāvśo to te rātre nahī̃ sūe.}
\end{quote}

\begin{quote}
If you frighten the child, he will not sleep tonight.
\end{quote}

Past, with \gu{તેણે} (he or she …): \gu{તેણે} \gu{નાના} \gu{ભાઈને} \gu{ડરાવ્યો}, \emph{teṇe nānā bhāīne ḍarāvyo}.

\begin{cousinweb}
\textbf{\gu{ડરાવવું}} (\emph{ḍarāvvũ}) is the causative of \textbf{\gu{ડરવું}} (\emph{ḍarvũ}), \textquotedblleft{}to be afraid\textquotedblright{}, from Sanskrit \emph{darati}, \textquotedblleft{}is frightened\textquotedblright{} (CDIAL 6190).
\end{cousinweb}

\begin{grammarlens}[title={What you've built}]
One more word for saying what happens if.
\end{grammarlens}

\subsection*{Your turn}
\begin{itemize}
\raggedright
  \item \emph{Say it:} \emph{ḍarāvvũ}
  \item \emph{Say it:} \emph{ḍarāvvũ}, once more
  \item \emph{Say it:} \emph{jo tame bāḷakne ḍarāvśo to te rātre nahī̃ sūe.} — If you frighten the child, he will not sleep tonight.
  \item \emph{Recall:} say the Gujarati for peace of mind, then the Gujarati for a misgiving, then say \emph{ḍarāvvũ} again
\end{itemize}

\begin{tcolorbox}[breakable,colback=teal!4,colframe=teal!35!black,title={Before you move on}]
What does \gu{ડરાવવું} mean? (\textquotedblleft{}To frighten\textquotedblright{}.) Say it once more.
\end{tcolorbox}
